\subsection{Rule formalization}

M-EXRx treats the rule layer as the component that records the binding decision.
Each rule has the form
\[
r=\langle id, layer, priority, body, decision, obligations, prohibitions,
evidence, trace\rangle .
\]
The body follows the shape of a RuleML implication, and the consequent emits
decision, obligation, prohibition, and trace atoms. The decision vocabulary is
\texttt{allow}, \texttt{downgrade}, \texttt{ask\_clarification},
\texttt{repair\_required}, and \texttt{refuse}.

Conflicts are resolved by strict priority:
\[
\begin{array}{l}
red\_flags > injury/fatigue/environment > scope > evidence\\
> protocol > capacity > progression > user\_preference .
\end{array}
\]
Thus a preference for intervals cannot override chest pain, heat illness,
acute injury escalation, or missing prescription evidence.

For a candidate plan \(P\) and contract \(C\), an \texttt{answered} output is
legal only when \(P\) uses allowed actions, avoids forbidden actions, respects
volume and intensity budgets, supports prescriptive claims with eligible
evidence, and emits the required trace. Otherwise the system must downgrade,
ask for clarification, return a partial answer, or refuse. Evidence eligibility
is narrower than retrieval relevance: protocol rules and active action-library
entries may authorize prescription, while ordinary retrieved literature remains
explanation-only unless promoted into an approved rule or action.

\paragraph{Canonical decisions.}
Four rule patterns cover the benchmark interface:
\[
\begin{array}{ll}
\mathrm{Allow:} &
R1 \wedge allowed\_action \wedge eligible\_evidence
  \Rightarrow \texttt{allow}.\\
\mathrm{Downgrade:} &
fatigue\_signal \wedge requested\_hard\_session
  \Rightarrow \texttt{downgrade}.\\
\mathrm{Refuse:} &
medical\_red\_flag \wedge asks\_for\_workout
  \Rightarrow \texttt{refuse}.\\
\mathrm{Clarify:} &
missing\_profile \wedge individualized\_plan
  \Rightarrow \texttt{ask\_clarification}.
\end{array}
\]
The full rule listing is part of the supplementary artifact.
